\chapter{How to Take an Argument Apart}
\section{A Bronze Disk Unearthed from the Soil}
First, pick something up.

It is scarcely larger than a coin, made of bronze, with a short axle through its center, like a flattened abacus bead. Twenty-four centuries ago, Athenian craftspeople cast these in batches and issued them to citizens serving in court. It is a voting disk.

There are two kinds: a solid axle means not guilty, a hollow one guilty. A juror pinches the axle's ends between thumb and forefinger and drops the disk into a box. Others see the vote being cast, but not which kind. The Greeks designed this thoughtfully: the verdict must be public, but individual choices can remain secret. In essence, the device converts personal judgments into a collective verdict.

Today you can see them in the museum at Athens's ancient Agora, green with corrosion, lying as quietly as old buttons in glass cases. Yet these little things once decided life and death. The most famous occasion was in 399 BCE, when a seventy-year-old man defended himself before about five hundred Athenians holding bronze disks. His defense failed. Most disks dropped with hollow axles.

The man was Socrates, widely regarded as a founding figure of Western philosophy. His judges were not tyrants, but roughly five hundred ordinary Athenians, chosen by lot in the most democratic city-state of the time.

This chapter does not ask whether he should have been convicted. That is a historian's question. It asks something closer to you. Each juror believed they had reasons, but having reasons and having reasons that hold up are different things. When a view appears before you, from a teacher, an online post, a convinced friend, or your own fresh thought, how can you tell whether it stands?

You must learn to take it apart. Like a bicycle mechanic dismantling a wheel, recognize that a view is not one solid lump of iron. It is assembled from components. This chapter teaches disassembly.

\section{The Morning Five Hundred People Sat Together}
Now enter the scene with me.

Time: a day in 399 BCE. Place: an Athenian court. Not the solemn hall we might imagine, but something closer to an enclosed courtyard. Around five hundred jurors have arrived before dawn. They are not professional judges but ordinary citizens drawn by lot that morning: tanners, olive-oil sellers, vine growers, retired soldiers. Selection by lot itself expresses an Athenian principle: judgment should belong to ordinary people, not experts. They sit on benches, murmuring like a hive.

A water clock stands nearby, a large leaking vessel that measures speaking time. A fixed quantity drains for the prosecution, the same for the defense. Even time is shared in an Athenian court.

A young man named Meletus rises to prosecute. He reads three charges: failing to honor the city's recognized gods, introducing new divinities, and corrupting Athens's youth. Behind him stands the more influential Anytus, a practical man who made his fortune in tanning and is a pillar of the democracy.

Socrates stands as defendant. If you expect tears and pleas for mercy, you are mistaken. In Plato's later account, the old man rises, calmly teases the prosecutors' eloquence, then questions Meletus until the young man stammers. Laughter and jeers break from the benches. Heckling is ordinary business in Athenian courts.

The water runs out and the speeches end. Five hundred people pass two boxes, pinching the ends of their axles and dropping their bronze disks. Later accounts put the margin at only a few dozen votes. Had another thirty-odd people chosen the other box, the old man could have gone home for dinner.

A second ballot follows. Athenian rules require each side to propose a penalty after conviction; the jury chooses between them. The prosecutor asks for death. Ordinarily, Socrates should now propose exile: leave Athens, preserve life and dignity, and likely win the jury's agreement. He refuses. He rises and says he has lived with a clear conscience. The city should reward, not punish him, giving him free meals for life in the civic hall, like an Olympic champion. Then, apparently thinking the joke has gone too far, he adds a nominal fine as an alternative.

The benches erupt. The second ballot approves death by a larger majority than the conviction.

Remember this morning. The rest of the chapter answers the questions hidden in it.

\section{What Happened Before the Disk Fell?}
Let us lay out the conflict without rushing to an answer.

On one morning, the same jurors hear the same words and see the same people, then divide: some choose hollow axles, others solid. The difference is not the evidence, shared by everyone, but the invisible journey between evidence and conclusion.

A convenient explanation says those voting guilty were persuaded and the others were not. But what does persuaded mean? By reasons or fear? By an argument heard to the end or an old grievance recalled?

Five years before the trial, Athens had lived through a nightmare. In 404 BCE, it lost a long war to Sparta. Sparta installed a puppet regime of Athenians known as the Thirty Tyrants. They confiscated opponents' property and executed them in a reign of terror lasting more than half a year, until democratic forces overthrew them. Their leader Critias had, as a young man, been one of the pupils constantly following Socrates. So had Alcibiades, Athens's most gifted and controversial general, whose defection to Sparta nearly destroyed the city.

Now put yourself in a juror's shoes. You are an Athenian tanner who lost a brother in war and hid in a cellar under the tyrants. Before you stands Athens's most famous old man. The charge reads impiety, but your mind sees his two pupils who threw the city into turmoil. You hold the disk. What, between your fingers, is really deciding?

This court illuminates the chapter's central question:

\textbf{What conditions must an argument meet to hold up?}

Is a conclusion that happens to be right enough? A persuasive sound? A majority nodding? Or is an argument like a bridge, needing sound timber and every nail between evidence and conclusion properly placed before it can carry a person?

A more uncomfortable question follows. If even Athens's democratic court, with five hundred citizens seriously fulfilling their duty, could vote under the pressure of fear and grievances, why are we sure we do better when forwarding a message, raising a hand in class, or privately deciding what something means?

Before reading on, weigh it yourself. If you were one of the five hundred, knowing only what you now know, which box would receive your disk?

\section{Make the Case for the Hollow Axles}
People today easily take sides in this ancient trial: Socrates was wise, the Athenians a mob, story over. That is too quick. We will practice a skill that recurs throughout this book and deserves a name: \textbf{steelmanning}.

You may know the straw man: replace another person's view with a stupid version, then knock it down. Steelmanning does the opposite. Before answering someone, present their strongest reasons, including those they have not fully articulated, faithfully and perhaps more clearly than they did. Straw is soft and easy to strike. Steel is hard; overcoming it counts.

Let us state the hollow-axle voters' case as fully as possible.

\textbf{First: their fear did not come from nowhere.} The city's memory was fresh. The Thirty Tyrants had fallen only five years earlier, and every household remembered Critias. Alcibiades's betrayal had nearly ended Athens. Socrates need not have taught these two to do evil for a reasonable question to remain. Does a teacher who spends his days urging young people to question established authority bear no responsibility for pupils' extreme acts? Teachers need not answer for their students' entire lives, but is the atmosphere of contempt for tradition and admiration for argumentative skill they create entirely innocent? The question is still recognizable when a celebrated teacher's pupil gets into trouble today.

\textbf{Second: everyday resentment was real.} For years, Socrates had stopped self-satisfied, respectable people in the square, politicians, poets, craftspeople, and questioned them until their inconsistencies left them publicly embarrassed. They hated it, and young spectators' laughter deepened the hatred. How many jurors or their friends had suffered such humiliation? More importantly, many fathers watched sons neglect work and apprenticeships to follow Socrates learning to question. For them, corrupting the young was no abstract charge. It happened at their dinner tables.

\textbf{Third: the procedure was clean.} This is easiest to overlook. By Athenian law, the trial broke no rule: public accusation, jurors by lot, equal speeches, secret votes, two ballots. The democrats had even shown restraint. On returning to the city five years earlier, they issued an amnesty barring prosecution of old grievances. Anytus could not bring the political charge that Socrates had taught Critias, so used the lawful religious route. The hollow-axle voters were not a lynch mob. They followed rules in one of their world's most procedure-conscious institutions.

Now the other side. Socrates's supporters have a full case too. He had served the city as a soldier, both attacking and covering retreats. Under the tyrants, ordered to arrest an innocent citizen for execution, he disobeyed and went home, a fact not every guilty voter may have remembered. He never charged for lessons, unlike the sophists who made fortunes teaching debate. As for corrupting the young, his questions urged them to care about their souls rather than money and reputation.

Both cases are now before us. Notice the method: before deciding who is right, we extracted and displayed \textbf{the strongest part} of each argument. Steelmanning makes the question harder, not simpler, but that is its real shape. Cases that seem easy to judge often reflect corners you have cut.

\section{Thinking Bridge: An Argument's Four Parts}
Reasons are on the table. Now we need tools. Serious thought without tools often means spending longer with a prejudice. This tool is \textbf{argument dissection}. Any view trying to make you believe something can be divided into four parts.

\textbf{First, the conclusion.} The statement the speaker ultimately wants you to accept. Socrates is guilty; phones should be banned; this film is bad. The conclusion is the destination all other statements serve. A simple way to find it is to ask which single sentence the speaker would most want preserved if only one could remain.

\textbf{Second, the reason.} The grounds for belief. Reasons answer why the conclusion is true. They are not the conclusion itself. In ``He should be punished because he broke the law,'' punishment is the conclusion and breaking the law the reason. Many arguments stall because someone repeats a conclusion ten ways without one reason: ``This rule is unreasonable! Completely unreasonable! Ridiculous!'' Volume is not a reason.

\textbf{Third, the evidence.} Reasons also need support. Why believe he broke the law? Testimony, records, physical evidence, data. Evidence is the component touching the world directly and, in principle, open to outside inspection. Others can look, investigate, and check. That distinguishes it from a reason: reasons form the logical chain; evidence is the ground beneath it.

\textbf{Fourth, hidden premises.} This is the least visible component. Nearly every argument relies on unspoken judgments treated as obvious. Meletus's argument says: Socrates dishonors the city's gods, therefore he is guilty. Several premises lie between: dishonoring gods should count as crime; the city's list of gods is the right one; the observed behavior really counts as dishonor. None has been argued for, yet each bears weight. Remove one and the argument falls.

Hidden premises are unseen piles beneath a bridge. To find them, ask \textbf{what the speaker assumes in moving from reason to conclusion}. A classmate says, ``This film scores 9.2 and you haven't seen it? Go see it!'' The conclusion is that you should watch it, the reason its high score, the hidden premise that a highly rated film is worth your time too. That may not hold: you might hate the genre. Even one conversational sentence contains premises, let alone an essay or speech.

Draw the four parts and you have an \textbf{argument map}. Use simple arrows: evidence below supports reasons, reasons point to the conclusion, and hidden premises sit beside the arrows because the arrows depend on them. Next time a textbook passage, speech, or video script tries to persuade you, draw this map. An argument that cannot be mapped may not be clear even to its speaker. One that can be drawn but has an unsupported link shows where to question.

Practice on Meletus's central charge. Conclusion: Socrates corrupts youth. Reason: his followers learn to defy fathers and despise authority. Evidence: testimony from some fathers and young people. At least two hidden premises carry the argument. First, defying fathers and authority is always bad. Second, Socrates causes these changes, not something else, such as growing up, war's intergenerational trauma, or turmoil in the city itself. The second premise, never seriously checked in that court, decides life and death. Two millennia later, we recognize a named trap in it, which this chapter will explain.

\section{A Few Words from the Old Man}
Now read a short primary source. Socrates never wrote anything himself. His pupil Plato recorded, or reconstructed, his defense in the Apology, one of the most widely transmitted ancient Western texts. After the death sentence, when words could no longer change the outcome, the old man says:

\begin{quote}
``The unexamined life is not worth living.''
\end{quote}
---Plato, Apology 38a

One sentence, but take your new scalpel to it. Conclusion: an unexamined life is not worth living. Reason: examination, continually asking what one believes and why, is fundamental to life. Hidden premises: life's value lies not only in living and enjoying, but in being understood and examined; whether it is worth living can and should be discussed. Those premises still provoke disagreement. ``Why think so much? Just be happy'' challenges the hidden premise, not logic itself, but the unwritten bridge plank.

A further passage near the Apology's opening is worth paraphrasing. Socrates tells the jury that the prosecutors spoke so brilliantly and persuasively that even he nearly believed them. Yet almost nothing they said was true, whereas at his age he would speak as he usually did in the marketplace, following his thoughts. The point is that persuasion comes in two kinds: elegant rhetoric and the weight of facts, often traveling separately. An old man at the boundary of life and death warned five hundred listeners against pleasing words twenty-four centuries ago. The warning has not expired.

Turn back to China. At roughly the same time, the Warring States Mohists seriously studied how to make reasoning clear, calling the practice bian, disputation. A programmatic sentence in the Mozi's ``Lesser Selection'' says:

\begin{quote}
``Use names to refer to realities, statements to express intentions, and explanations to set out reasons.''
\end{quote}
---Mozi, ``Lesser Selection''

Roughly: names identify things, sentences convey meaning, arguments supply grounds. Notice the final word, gu: grounds or reasons. The Mohists put giving reasons at language's center. What you say is not enough; why it holds is what matters. They also systematically examined rules and traps of debate, such as why one person's being a robber does not imply that everyone where robbers live is a robber. Athens's square and China's schools, separated by a continent and unaware of one another, worked in the same century on the same question: how can words hold up? This is no coincidence. Once people seriously discuss public affairs, the craft of dismantling views is bound to emerge.

\section{One Court, Four Kinds of Reasoning}
With a scalpel in hand, we can address a deeper question. How many routes lead from evidence to conclusion?

There is more than one. Recognizing the routes and their different reliability is this chapter's most practical lesson. Return to the Athenian court and use the same case to try four in turn.

\textbf{First, deduction: from general to particular, preserving truth at every step.}

Deduction states a general rule, places a particular case under it, and lets the conclusion follow. Not long after Socrates's death, his pupil's pupil Aristotle studied this thoroughly in works later collected as the Organon, meaning instrument. He saw logic as the blade every field must sharpen before beginning. A standard example has circulated in logic classes for more than two thousand years:

\begin{quote}
All humans are mortal. Socrates is human. Therefore Socrates is mortal.
\end{quote}
This reasoning has an extraordinary property: if both premises are true, the conclusion \textbf{cannot} be false. No further investigation is needed; the premises already contain the conclusion, and deduction releases it. Meletus attempts this form: the city's laws punish impiety, Socrates is impious, therefore he should be punished. Perfect form, but perfect form guarantees only that true premises yield a true conclusion. Deduction cannot establish whether the premises themselves are true. That limitation is crucial, and the next section examines it.

\textbf{Second, induction: from many particulars to a general claim, offering strength but no guarantee.}

Induction observes one case after another, then generalizes. The charge of corrupting youth is essentially inductive: one follower defies his father, another neglects work, a third does something similar, so Socrates corrupts youth.

Induction is a mainstay of science and ordinary life. Tomorrow's sunrise is an inductive expectation. But one rule is ironclad: \textbf{the conclusion is never locked in}. Ten thousand white swans cannot guarantee the next is not black. Australia's black swans overturned Europe's assumption that all swans were white. Inductive conclusions therefore have degrees of strength. More numerous, varied cases that withstand counterexamples make a stronger conclusion. Those three wayward youths? Too small a sample, and probably selected. Fathers who hate Socrates will not bring forward a young man who became more responsible after talking to him. Induction's most insidious trap is often not the reasoning but secretly selected cases.

\textbf{Third, analogy: from one case to another, on a borrowed bridge with limited strength.}

Analogy says A and B resemble each other in some respects; A has a property, so B probably does too. The court's finest analogy is Socrates's own. He asks Meletus: You say all Athens improves young people and I alone corrupt them. Consider horses. Can any passerby train them well, or only a few skilled trainers? Clearly the latter. Why, then, can everyone educate youth except me? This paraphrases the analogy in the Apology.

Witty and powerful, it wins Socrates points. But analogy borrows a bridge. Horses and young people may both require skilled training, yet differ elsewhere. Horses cannot freely choose teachers; young people can. Horse training's goal is clear; education's goals are disputed. Analogy properly offers insight, a new way to look. Misused as evidence, it claims similarity establishes identity. Arguments beginning ``Think of this class as a family'' or ``The state is like a ship'' deserve an extra question: alike, yes, but where are they unlike?

\textbf{Fourth, inference to the best explanation: finding the most economical account of facts.}

The first three move from premises toward conclusions. Another works backward: begin with a fact and ask which explanation best accounts for it. Detectives, doctors, and people diagnosing broken phones use it. Let us apply it to the chapter's greatest unresolved question: \textbf{why did most jurors choose hollow axles?}

At least three explanations compete for the same fact.

\textbf{Explanation one: they genuinely believed the charges.} Religion lay at the heart of Athenian public life. Impiety genuinely threatened the city because divine anger could bring disaster on everyone. For Athenians, impiety resembled a threat to public safety today. Jurors heard both sides and voted on evidence. Its case: their religious feelings were real, and Meletus's charge had standing in contemporary law. Its limit: why prosecute only in 399 BCE, after decades of Socratic questioning? Why not sooner?

\textbf{Explanation two: political reckoning wore religious clothes.} The true triggers were Critias, Alcibiades, and fear left by defeat and tyranny. Amnesty blocked political prosecution; religious charges provided the only legal route. Its case: it explains both the timing, a few years after democratic restoration, and the involvement of a practical democrat such as Anytus. Its limit: it makes five hundred jurors too much like puppets. Their votes were secret and uncoerced. If most truly saw a harmless old man, political resentment might not bridge the margin of a few dozen votes.

\textbf{Explanation three: Socrates lost his own case.} Conviction had a narrow margin. During sentencing, the conventional proposal of exile could have saved him, but demanding free meals pushed the middle against him. Many ancient and modern readers think he meant this death to complete his life. Its case: it neatly explains the larger majority in the second ballot. Its limit: it cannot explain the first. Before that ballot, he had done none of this, yet hundreds still voted guilty.

How do we choose? Inference to the best explanation uses plain tests. Which explains more details? Which needs fewer coincidences? Which better fits known human behavior and institutions? Weighing these, the first appears too naive, the third explains only half, and the second is strongest but erases genuine jurors' judgments. The honest conclusion is that \textbf{each captures part of the truth, and historians still have no single agreed verdict}. This is not evasive compromise, but marking the evidence's boundary. Admitting that available material takes us only this far is itself part of reasoning skill.

Consider how other civilizations sharpened the blade. In India, repeated debates among Buddhist and Brahmanical schools produced a finely developed study of inference, hetuvidya. Its standard argument has three members: thesis, reason, example. A classic goes roughly: sound is impermanent, because it is produced; produced things are impermanent, as a pot is. Notice the example. The reasoner must supply not only a rule but a testable instance, stitching deduction and induction together and guarding against empty debate. Centuries later, scholars in the Arabic-speaking world, including Ibn Sina, known in Latin as Avicenna, organized and advanced logic transmitted from Greece. Their works eventually reached Europe and became standard medieval university reading. Dissecting arguments is a craft independently invented and mutually refined by human beings. It belongs to no single civilization, which is part of why it deserves trust.

\section{Further Challenge: Right Answer, Wrong Route}
Now for this chapter's weightiest theory, from logic. Its core is a distinction we must keep clear: \textbf{validity} and \textbf{truth}.

Calling reasoning valid concerns its \textbf{form}: true premises cannot lead to a false conclusion. Calling its premises true concerns its \textbf{material}. These are separate checkpoints, and logic's most important lesson is that \textbf{neither guarantees the other}.

First test why validity does not guarantee true premises:

\begin{quote}
All birds can fly. Penguins are birds. Therefore penguins can fly.
\end{quote}
The form is perfect, identical to the mortal Socrates example, and entirely valid. The conclusion is false. The error lies in the first premise: not all birds fly. This is valid but unsound reasoning. The machine works, but bad material goes in and a bad product comes out. It teaches a lifelong discipline: \textbf{before admiring watertight reasoning, go back and check its premises}. The more elaborate and rigorous an argument, the more easily its elegance distracts us from unchecked material at the door.

Now the reverse, an easily misjudged counterexample: \textbf{a true conclusion does not guarantee correct reasoning}. Test this:

\begin{quote}
All mammals are animals. Cats are animals. Therefore cats are mammals.
\end{quote}
The conclusion is true: cats are mammals. But the reasoning is \textbf{wrong}. How can we prove it? Keep the form and change the material:

\begin{quote}
All mammals are animals. Fish are animals. Therefore fish are mammals.
\end{quote}
Same form, failed conclusion. The form is unsafe; the earlier success was luck. Remember this example for life, because it is common. Someone reaches a correct conclusion through broken logic, gains confidence in that logic, and uses it again. Next time luck may fail. \textbf{The world judges conclusions; form judges reasoning. These are separate accounts.}

Bring the distinction back to Athens and much becomes clearer. Meletus's argument has valid form but fails at its premises: whether Socrates's conduct counts as impiety was never seriously verified. Some defenses of Socrates may have moving conclusions but forms that fail when their contents change. What the ancient court lacked was these two separate checks.

Another subtler trap hides in one word. The Warring States thinker Gongsun Long left a ``Discourse on the White Horse,'' whose central claim is four Chinese characters:

\begin{quote}
``A white horse is not a horse.''
\end{quote}
---Gongsun Longzi, ``Discourse on the White Horse''

It sounds like mere provocation. How can a white horse not be a horse? Yet Gongsun Long argues fluently: horse names a shape, white names a color, and white horse combines shape and color, so white horse is not identical to horse. The audience becomes entangled because the meaning of is quietly changes. If not means the concepts are not identical, the statement is true: the expressions differ in meaning. If it means a white horse does not belong to the class of horses, it is false. A word carries two meanings, the speaker slides between them, and the reasoning collapses. Philosophers later called this ambiguity, and named the defense too: fix meanings before reasoning. Next time a verbal knot confuses you, do not first admire it. Ask whether each word keeps the same meaning throughout.

Finally, mark the theory's limit. Logic checks form, questions premises' sources, and fixes meanings, but cannot verify facts for you or choose what is worth pursuing. Two people can share correct logic yet reach different conclusions because they care about different things. Logic is a good knife; it does not choose what to cut. In everyday quarrels, as the next section shows, simply recognizing faulty form already avoids many traps.

\section{The Class Chat: Six Common Traps}
This twenty-four-century-old question goes to court in your pocket every day.

Suppose your year group is debating a total ban on students bringing phones to school. Hundreds of messages fill the class chat. Watch a short exchange, then dissect it. Every move from Athens's square is still alive, only the screen has changed.

\textbf{Trap one: the personal attack.}

A student writes, ``The administrator proposing this can't even use a smartphone. Why discuss anything he says?''

Dissect it. Whether the ban is good depends on reasons and evidence. Whether its proposer can use a phone is another matter. Attacking the person instead of the view is a personal attack, or ad hominem. Its illusion is that discrediting a person also defeats the idea. In fact, the idea remains untouched.

Beware an easy misjudgment, though: not every person-directed statement is an ad hominem. ``This reviewer was paid by the phone manufacturer, so we should be cautious about their specifications'' checks an evidence source, exactly the skill this chapter teaches. Attacking character to avoid an argument is fallacious; examining financial interests to test evidence is legitimate. The dividing question is whether the speaker is examining the view or evading it.

\textbf{Trap two: the straw man.}

A teacher proposes, ``During classes, put phones in storage pockets at the front and collect them after school.'' Someone immediately replies, ``You want to cut students off from their families! What about emergencies? Are we making school a prison?''

The original proposal of classroom storage has been replaced with total isolation and imprisonment, a weaker, more absurd version, then attacked. This is the straw man: build an easy target, knock it down, declare victory. The defense is simple and courteous. Before rebutting, restate the other person's view and ask whether you understood it correctly. That is the discipline of steelmanning: build steel, not straw.

\textbf{Trap three: the false dilemma.}

``Either ban phones completely or lose all classroom control. Which do you choose?''

Are these really the only possibilities? Classroom storage, time limits, dormitory-only restrictions, rules by year group: many arrangements lie between. A false dilemma folds a whole field of possibilities into two extremes and forces a choice between bad options. Ask why there are only two. Usually the dilemma dissolves.

\textbf{Trap four: appeal to the majority.}

``Let's vote. Eighty percent of the year favors a ban. What can the minority still say?''

There are two levels here. If the dispute concerns the school life we want, it is a matter of preference and choice. Voting is perfectly legitimate; majority decisions select class representatives and uniforms. If the question is whether phones lower grades, it is factual and cannot be settled by raised hands. Geocentrism was once universal consensus, a hundred percent of the votes, yet Earth still orbited the sun. Truth is not distributed by head count. The mistake is not voting, but \textbf{using votes to decide facts}.

\textbf{Trap five: hasty generalization.}

``My cousin's school allowed phones for a year and the whole year's grades collapsed. Phones kill achievement.''

One school, one year group, one year: applying that sample to tens of millions of students is like testing a lake with one cup of water. Remember induction's rule: number and diversity of cases determine a conclusion's strength. Someone's cousin is one case, not a chain of evidence. When you hear ``Someone I know...,'' politely ask what followed, how many others there were, and what about the cases not told.

\textbf{Trap six: reversed causation.}

``The data are clear. Students who spend longer on phones have worse grades. Phones harm students. Case closed.''

Long phone use and poor grades often occur together. That is correlation. Between occurring together and one causing the other lies an unargued journey. At least three possibilities exist: phone use lowers grades, A causes B; struggling students seek achievement on phones, B causes A; or both arise from a third cause, poor sleep, family upheaval, lost interest in study, C causes both. Ice-cream sales and drownings rise and fall together too, not because ice cream causes drowning but because summer increases both. Distinguishing these possibilities is basic to reading case-closed news. Athenians who concluded that Socrates caused young people's deterioration fell into this trap: they never checked the third possibility.

All six traps are marked. Look back at the chat: personal attacks change accounts, straw men are readily made, false dilemmas sound commanding, majority appeals earn applause, generalizations tell vivid stories, and reversed causation loves wearing data's clothes. These moves survive not because they are correct, but because they \textbf{save effort}. Dismantling a real view is difficult; making a straw man takes three seconds. Recognizing fallacies is not a weapon for always winning quarrels. It helps you distinguish a bridge from a pit when you are ready to be persuaded.

\section{Over to You: When Should Dissection Stop?}
Return to the five hundred bronze disks.

You now know that arguments divide into conclusions, reasons, evidence, and hidden premises. Deduction, induction, analogy, and best explanation offer routes of different strength. Valid form does not guarantee true material, and a true conclusion does not guarantee a sound route. Six effort-saving traps wait along those routes. Had these tools stood outside the Athenian court, would the vote have changed? Perhaps, perhaps not. Tools work only in willing hands.

This brings us to a final question without a standard answer.

Steelmanning asks you to give the strongest version of another person's reasons before replying. It is good craft in class chats, courts, and arguments with yourself. But consider an uncomfortable case: what if the view's central premise is that some category of people is naturally inferior? Must you make that case as strong as possible too?

One side has strong reasons. Without taking it apart, you do not know where it fails. Declaring a view unworthy of discussion wins the scene but loses the reasoning. When it returns tomorrow in different clothing, you will not recognize it. And who controls the list of views unworthy of examination? Athenians in 399 BCE put teachers questioning tradition on a list of people not worth taking seriously.

The other side also has strong reasons. Restating a harmful view clearly, forcefully, and elegantly helps sharpen its weapon. Some claims half-win simply by gaining the status of worth serious discussion. Attention is a resource, and a limited one. Spending it on every view that knocks can itself be irresponsible.

Where is the line? Which views deserve your full toolkit, and which merit only ``I reject that premise''? Who draws the line, and on what grounds?

There is no standard answer. But take the question into your next quarrel, trending topic, or essay that makes you nod repeatedly. Since the morning bronze disks fell into those boxes, humanity has practiced one craft: before casting your vote, take apart the view before you and look inside.
